\section{Model Architecture}
\label{sec:model}

Figure~\ref{fig:overall} presents the high-level design of \tool{}'s approach.
We model the problem of throughput estimation as a regression problem: given the assembly input, \tool{} predicts the throughput of the instruction sequence as a real-valued number.
At the core of \tool{} is a hierarchical multiscale RNN~\cite{dag-rnn1, dag-rnn2} that sequentially processes all instructions in the basic block and outputs an embedding, which \tool{} then uses to directly estimate the throughput.
Altogether, we decompose the end-to-end model into the following stages: \emph{canonicalization}, \emph{embedding} and \emph{estimation}.

\subsection{Canonicalization}
\label{sec:canonicalization}

The canonicalization stage converts the assembly input into a more structured form, dictated by the syntax of the assembly instructions.
\tool{} takes a compiled assembly block, disassembles it, and maps it to a list of instructions.
Each instruction consists of a list of tokens representing its operation code (opcode, e.g. \texttt{add}), source operands, and destination operands, separated by distinguished delimiter tokens.

For example, consider the instruction \texttt{mul ecx}, which multiplies the value in register \texttt{ecx} with \texttt{eax}, and places the result into registers \texttt{edx} and \texttt{eax}.
Note that the source operand \texttt{eax} and both of the destination operands \texttt{eax} and \texttt{edx} are implicit in the Intel syntax \texttt{mul ecx}.
The final canonicalized set of tokens for the instruction is:
$$(\texttt{mul},\,\textit{<S>},\,\texttt{eax},\,\texttt{ecx},\,\textit{<D>},\,\texttt{edx},\,\texttt{eax},\,\textit{<E>})$$
where the bracketed tokens are the delimiters representing the break between the opcode, source, and destination operands.

Assembly code permits more than just register operands, such as constants and memory operands.
We map all constants (e.g. integer constants, memory addresses, etc.) to a single \texttt{CONST} token.
We demarcate memory operands (consisting of a base address, and an optional offset and displacement) by surrounding them with \texttt{<M>} and \texttt{</M>} delimiter tokens.
We present the full canonicalization scheme in Appendix~\ref{app:token-grammar}.

\subsection{Embedding}
\label{sec:nnembeddings}

{\tool}'s embedding stage takes a canonicalized token stream of instructions, and for each instruction produces an \emph{embedding}: a representation of an instruction as a real-valued vector in a high-dimensional space.
The first step is the \textit{token layer}, which maps a given token to an embedding.
We implement the token layer by mapping each token in the sequence to an $n$-dimensional vector by learning a linear transformation of the one-hot token vectors (this is equivalent to learning a lookup table).

\tool{} then maps the sequence of token embeddings to an embedding for each instruction in the basic block.
We call this the \textit{instruction layer}.
Because each instruction can have a variable number of tokens depending on its number of source and destination operands, the size of the input to the embedding stage is variable.
We therefore implement the instruction layer with a sequential Recurrent Neural Network (RNN) architecture with Long Short Term Memory (LSTM)~\cite{lstm} cells.

Figure~\ref{fig:overall} presents the operation of our RNN-based instruction embedding approach on a small example.
The bottommost row shows the original assembly input.
The second row shows the sequence of tokens for each instruction.
The third row (the token layer) shows the sequence of {\em token embeddings}, e.g. $v_\texttt{mov}$, which are mapped directly from each syntactic token.
The fourth row (the instruction layer) shows the application of an LSTM to reduce the token embedding sequence into the final instruction embedding, $h_{\texttt{mov}}$.

\subsection{Prediction}
\label{sec:nnoverview}

The final prediction comes from the \textit{prediction layer}, which maps a basic block (a sequence of instruction embeddings) to a throughput value.
This is again implemented with an RNN with LSTM cells, which has entirely disjoint weights from the LSTM in the instruction layer.
This corresponds to the topmost layer in Figure~\ref{fig:overall}.
Using the final output from the instruction LSTM ($h_{\texttt{block}}$), \tool{} predicts the basic block's throughput with a linear layer.
Specifically, {\tool} computes $w \cdot h_{\texttt{block}} + b$, where $w$ is a learned weight vector and $b$ is a bias.
This produces a final real-valued number that represents the network's throughput prediction.

The hierarchical combination of the RNN in the instruction layer and the RNN in the prediction layer has several benefits over a non-hierarchical model:
\begin{itemize}
\item Memory and backpropagation paths are significantly shorter than a model using a non-hierarchical (i.e., single) RNN.
  The average length of a block in our dataset is $6.04$ instructions, and the average length of an instruction is $7.97$ tokens.
  The average length of a token-level RNN across the entire basic block would instead be about $48$ RNN cell applications, more than three times as long as a path through the hierarchical RNN.
\item
  Instructions are embedded atomically: the prediction layer is only able to generate throughput estimates at specific points in the overall token stream, i.e. after complete instructions.
  This means that the network does not have an obligation to produce states that correspond to predictions at points in between instructions.
\end{itemize}
We compare this hierarchical architecture against other architecture choices in Section~\ref{sec:exploration}, showing the efficacy of \tool{}'s hierarchical architecture.
